\documentclass[aspectratio=169,11pt]{beamer}
\usepackage[utf8]{inputenc}
\usepackage[T5]{fontenc}
\usepackage[vietnamese]{babel}
\usepackage{lmodern}
\usepackage{graphicx}
\usepackage{booktabs,tabularx,array}
\usepackage{tikz}
\usepackage{pgfplots}
\pgfplotsset{compat=1.18}
\usetikzlibrary{positioning,arrows.meta,calc}
\usetheme{Madrid}
\usecolortheme{spruce}
\setbeamertemplate{navigation symbols}{}
\setbeamertemplate{section in toc}[sections numbered]
\setbeamercolor{section in toc}{fg=red!70!black}
\setbeamercolor{subsection in toc}{fg=blue!70!white}
\setbeamercolor{section number projected}{bg=orange,fg=white}
\date{Tháng 11, 2026}
\author{Nhóm Propensity VSF}
\institute{EDA dữ liệu Dunnhumby}
\title{Hiểu dữ liệu giao dịch Dunnhumby}
\subtitle{Cấu trúc bảng, raw feature và tín hiệu hành vi}
\titlegraphic{\IfFileExists{fit-logo.jpg}{\includegraphics[width=3.2cm]{fit-logo.jpg}}{}}
\setbeamertemplate{title page}{%
\centering\vspace{0.5cm}\inserttitlegraphic\par\vspace{0.25cm}
\begin{beamercolorbox}[sep=8pt,center,rounded=true,shadow=true]{title}\usebeamerfont{title}\textbf{\inserttitle}\par\end{beamercolorbox}
\vspace{0.25cm}{\usebeamerfont{subtitle}\insertsubtitle\par}\vspace{0.35cm}
\begin{beamercolorbox}[sep=5pt,center]{author}\insertauthor\end{beamercolorbox}
\begin{beamercolorbox}[sep=5pt,center]{institute}\insertinstitute\end{beamercolorbox}
\vspace{0.25cm}{\small\insertdate\par}}
\newcommand{\src}[1]{\par\vfill{\tiny\color{gray}Nguồn: #1}}
\newcommand{\stat}[2]{{\Large\bfseries\color{structure.fg}#1}\quad {\small #2}\par\medskip}
\newcommand{\code}[1]{\texttt{\detokenize{#1}}}
\newcolumntype{Y}{>{\raggedright\arraybackslash}X}
\newcolumntype{C}[1]{>{\centering\arraybackslash}p{#1}}

\begin{document}

\begin{frame}[plain]\titlepage\end{frame}

\begin{frame}{Mục lục}
\tableofcontents
\end{frame}

\section{Bảng giao dịch đang chứa dữ liệu gì}

\begin{frame}{Bảng transaction ghi nhận product line trong basket}
\begin{columns}[T]
\column{.47\textwidth}
\stat{2.595.732}{product lines}
\stat{276.484}{baskets}
\stat{2.500}{households}
\stat{711}{ngày quan sát}
\column{.49\textwidth}
\textbf{Một dòng mô tả}
\begin{itemize}
\item Một sản phẩm trong một basket mua hàng.
\item Household, cửa hàng và thời điểm gắn với basket.
\item Số lượng, doanh số và ba khoản discount gắn với product line.
\end{itemize}
\medskip
\textbf{Phạm vi}
\begin{itemize}
\item 92.339 sản phẩm và 582 cửa hàng.
\item 102 tuần, tương đương khoảng 23 tháng lịch sử tương đối.
\item Không có null trong 12 raw columns.
\end{itemize}
\end{columns}
\src{Tính trực tiếp từ data/raw/Dunnhumby/transaction\_data.csv.}
\end{frame}

\begin{frame}{Grain được xác nhận bằng dữ liệu}
\begin{columns}[T]
\column{.54\textwidth}
\centering
\begin{tikzpicture}[
node distance=8mm,
box/.style={draw=structure.fg,rounded corners,fill=green!8,align=center,minimum width=3.5cm,minimum height=.8cm},
line/.style={-{Latex[length=2mm]},thick,draw=structure.fg}]
\node[box] (hh) {\code{household_key}\\entity};
\node[box,below=of hh] (basket) {\code{BASKET_ID}\\basket event};
\node[box,below=of basket] (line) {\code{BASKET_ID + PRODUCT_ID}\\unique product line};
\draw[line] (hh)--node[right,font=\scriptsize]{1:N baskets}(basket);
\draw[line] (basket)--node[right,font=\scriptsize]{1:N product lines}(line);
\end{tikzpicture}
\column{.42\textwidth}
\textbf{Kiểm chứng}
\begin{itemize}
\item 2.595.732 cặp \code{BASKET_ID + PRODUCT_ID} khác nhau.
\item Không có cặp grain trùng.
\item Một basket chỉ thuộc một household, một ngày, một store và một thời điểm.
\item Median 5 product lines mỗi basket, p90 là 24.
\end{itemize}
\medskip
\alert{Kết luận:} grain phù hợp để tổng hợp lên basket, ngày và household mà không nhập nhằng khóa.
\end{columns}
\src{Group-by và uniqueness checks trên transaction\_data.csv.}
\end{frame}

\begin{frame}{DAY là chỉ số ngày tương đối, không phải ngày lịch}
\begin{columns}[T]
\column{.55\textwidth}
\centering
\begin{tikzpicture}[x=.82cm,y=1cm]
\fill[green!18] (0,0) rectangle (5,1);
\fill[orange!25] (5,0) rectangle (12,1);
\draw[thick] (5,-.15)--(5,1.15);
\node[align=center] at (2.5,.5){\textbf{WEEK\_NO = 1}\\DAY 1--5};
\node[align=center] at (8.5,.5){\textbf{WEEK\_NO = 2}\\DAY 6--12};
\node[below,font=\scriptsize] at (0,0){DAY 1};
\node[below,font=\scriptsize] at (12,0){DAY 12};
\end{tikzpicture}
\vspace{.45cm}
\[
\texttt{WEEK\_NO}=\left\lfloor\frac{\texttt{DAY}+1}{7}\right\rfloor+1
\]
\column{.41\textwidth}
\textbf{Điều dữ liệu xác nhận}
\begin{itemize}
\item DAY chạy liên tục từ 1 đến 711.
\item WEEK\_NO chạy liên tục từ 1 đến 102.
\item Công thức trên khớp toàn bộ 2.595.732 dòng.
\item Tuần đầu chỉ có 5 ngày vì calendar week bắt đầu trước DAY 1.
\end{itemize}
\medskip
\alert{Còn thiếu:} mốc ngày lịch tương ứng với DAY 1.
\end{columns}
\src{Quan sát ánh xạ DAY--WEEK\_NO trong transaction\_data.csv.}
\end{frame}

\begin{frame}{Transaction là bảng sự kiện trung tâm}
\centering
\begin{tikzpicture}[
node distance=7mm and 14mm,
box/.style={draw=structure.fg,rounded corners,fill=green!7,align=center,minimum width=3.1cm,minimum height=.7cm},
core/.style={draw=orange!80!black,very thick,rounded corners,fill=orange!12,align=center,minimum width=3.5cm,minimum height=1cm},
line/.style={-{Latex[length=2mm]},thick}]
\node[core] (t) {\code{transaction_data}\\product-line events};
\node[box,above left=of t] (p) {\code{product}\\PRODUCT\_ID};
\node[box,above right=of t] (h) {\code{hh_demographic}\\household\_key};
\node[box,left=of t] (c) {\code{causal_data}\\PRODUCT\_ID, STORE\_ID, WEEK\_NO};
\node[box,right=of t] (ct) {\code{campaign_table}\\household\_key};
\node[box,below left=of t] (cp) {\code{coupon}\\PRODUCT\_ID};
\node[box,below right=of t] (cr) {\code{coupon_redempt}\\household\_key, DAY};
\draw[line] (t)--(p); \draw[line] (t)--(h); \draw[line] (t)--(c);
\draw[line] (t)--(ct); \draw[line] (t)--(cp); \draw[line] (t)--(cr);
\end{tikzpicture}
\vspace{.25cm}
\begin{columns}[T]
\column{.48\textwidth}\small
Toàn bộ 92.339 PRODUCT\_ID phát sinh giao dịch đều có trong product master.
\column{.48\textwidth}\small
Demographic chỉ phủ 801/2.500 households, nên phải giữ nhóm ``unknown demographic''.
\end{columns}
\src{Schema và key coverage của các CSV trong data/raw/Dunnhumby.}
\end{frame}

\section{Từng raw feature là gì}

\begin{frame}{Mười hai cột tạo thành bốn nhóm feature}
\small
\begin{tabularx}{\textwidth}{@{}p{2.5cm}Y Y@{}}
\toprule
\textbf{Nhóm} & \textbf{Raw columns} & \textbf{Vai trò trong phân tích} \\
\midrule
Entity và key & \code{household_key}, \code{BASKET_ID}, \code{PRODUCT_ID}, \code{STORE_ID} & Nhóm dữ liệu theo khách hàng, basket, sản phẩm và điểm bán \\
\addlinespace
Temporal & \code{DAY}, \code{WEEK_NO}, \code{TRANS_TIME} & Xác định thứ tự giao dịch, recency, cadence và khung giờ \\
\addlinespace
Purchase measures & \code{QUANTITY}, \code{SALES_VALUE} & Đo lượng mua và giá trị mua \\
\addlinespace
Promotion measures & \code{RETAIL_DISC}, \code{COUPON_DISC}, \code{COUPON_MATCH_DISC} & Đo mức giảm giá và phản ứng với coupon \\
\bottomrule
\end{tabularx}
\medskip
\begin{block}{Nguyên tắc}
Raw column cung cấp bằng chứng. Feature mô hình phải được tổng hợp trong observation window và chỉ dùng thông tin có trước snapshot.
\end{block}
\src{Feature catalog của notebook 01 và transaction schema.}
\end{frame}

\begin{frame}{Entity và key features}
\scriptsize
\begin{tabularx}{\textwidth}{@{}p{2.2cm}p{1.5cm}p{1.8cm}Y Y@{}}
\toprule
\textbf{Cột} & \textbf{Role} & \textbf{Cardinality} & \textbf{Nghĩa nghiệp vụ} & \textbf{Tín hiệu có thể tạo} \\
\midrule
\code{household_key} & Entity ID & 2.500 & Mã household ẩn danh & Khóa tạo lịch sử, RFM và nhãn propensity \\
\addlinespace
\code{BASKET_ID} & Event ID & 276.484 & Mã basket mua hàng & Frequency, basket size, basket value \\
\addlinespace
\code{PRODUCT_ID} & Product ID & 92.339 & Mã sản phẩm trên product line & Preference khi join product taxonomy \\
\addlinespace
\code{STORE_ID} & Store ID & 582 & Mã cửa hàng phát sinh basket & Store affinity, diversity và channel proxy \\
\bottomrule
\end{tabularx}
\medskip
\begin{columns}[T]
\column{.48\textwidth}
\textbf{Household lặp lại đủ dài}
\begin{itemize}\scriptsize
\item Median 79 baskets và 70 active days.
\item Median span 627 ngày.
\item 2.461/2.500 households có span ít nhất 365 ngày.
\end{itemize}
\column{.48\textwidth}
\textbf{Lưu ý mô hình}
\begin{itemize}\scriptsize
\item Không scale ID như biến số.
\item PRODUCT\_ID và STORE\_ID cần aggregate hoặc encode theo lịch sử.
\item Chia train/test theo thời gian, không random theo dòng.
\end{itemize}
\end{columns}
\src{Cardinality và household-level aggregation từ transaction\_data.csv.}
\end{frame}

\begin{frame}{QUANTITY và SALES\_VALUE đo hành vi mua}
\begin{columns}[T]
\column{.48\textwidth}
\textbf{QUANTITY}
\begin{itemize}
\item Median = 1, p99 = 10.
\item 14.466 dòng bằng 0, chiếm 0,56\%.
\item Max = 89.638 và skewness = 13,64.
\item Đơn vị chưa được data dictionary cục bộ xác nhận.
\end{itemize}
\medskip
\textbf{Tín hiệu:} số lượng mua, bulk buying, basket units, category volume.
\column{.48\textwidth}
\textbf{SALES\_VALUE}
\begin{itemize}
\item Median = 2, p99 = 20.
\item 18.850 dòng bằng 0, chiếm 0,73\%.
\item Max = 840 và skewness = 15,91.
\item Tổng quan sát = 8,06 triệu đơn vị tiền tệ.
\end{itemize}
\medskip
\textbf{Tín hiệu:} monetary value, average basket value, spend trend.
\end{columns}
\begin{alertblock}{Điểm cần chốt}
IQR của QUANTITY bằng 0 vì Q1 = median = Q3 = 1. Không nên tự động loại mọi giá trị khác 1 như outlier.
\end{alertblock}
\src{Numeric profile trên 2.595.732 product lines.}
\end{frame}

\begin{frame}{Ba discount columns đo các cơ chế khuyến mãi khác nhau}
\scriptsize
\begin{tabularx}{\textwidth}{@{}p{2.4cm}C{1.4cm}C{1.4cm}C{1.4cm}Y@{}}
\toprule
\textbf{Cột} & \textbf{Min} & \textbf{Zero rate} & \textbf{Dòng âm} & \textbf{Tín hiệu hành vi} \\
\midrule
\code{RETAIL_DISC} & -180 & 49,80\% & 1.303.026 & Mức phụ thuộc vào giảm giá tại quầy, tỷ trọng chi tiêu có markdown \\
\addlinespace
\code{COUPON_DISC} & -55,93 & 98,60\% & 36.422 & Tần suất và giá trị sử dụng coupon \\
\addlinespace
\code{COUPON_MATCH_DISC} & -7,70 & 99,33\% & 17.449 & Phản ứng với chương trình match coupon \\
\bottomrule
\end{tabularx}
\medskip
\begin{columns}[T]
\column{.48\textwidth}
\textbf{Cách đọc}
\begin{itemize}\scriptsize
\item Giá trị âm nhiều khả năng là khoản giảm trừ khỏi sales.
\item Coupon columns zero-inflated vì phần lớn product lines không dùng coupon.
\end{itemize}
\column{.48\textwidth}
\textbf{Cần xác nhận}
\begin{itemize}\scriptsize
\item Đơn vị tiền tệ và quy ước dấu.
\item 36 dòng RETAIL\_DISC dương.
\item SALES\_VALUE đã là net hay gross sales.
\end{itemize}
\end{columns}
\src{Domain và numeric profiles từ transaction\_data.csv.}
\end{frame}

\begin{frame}{Ba biến thời gian phục vụ ba câu hỏi khác nhau}
\small
\begin{tabularx}{\textwidth}{@{}p{2.2cm}p{2.2cm}p{2.7cm}Y@{}}
\toprule
\textbf{Cột} & \textbf{Encoding} & \textbf{Miền quan sát} & \textbf{Tín hiệu có thể tạo} \\
\midrule
\code{DAY} & Ordinal day index & 1--711, đủ 711 ngày & Recency, active days, interpurchase gap, tenure \\
\addlinespace
\code{WEEK_NO} & Ordinal week index & 1--102, đủ 102 tuần & Weekly cadence, seasonality, rolling windows \\
\addlinespace
\code{TRANS_TIME} & Integer HHMM & 0000--2359, đủ 1.440 phút & Morning/evening affinity, shopping-time regularity \\
\bottomrule
\end{tabularx}
\medskip
\begin{block}{Vì sao TRANS\_TIME cần route riêng}
0930 và 1000 cách nhau 30 phút, không phải 70 đơn vị. Clock-time phải được parse thành phút trong ngày trước khi tính khoảng cách hoặc cyclic encoding.
\end{block}
\src{Temporal profile và kiểm tra HHMM từ transaction\_data.csv.}
\end{frame}

\begin{frame}{Các điểm dữ liệu cần xử lý bằng business rule}
\begin{columns}[T]
\column{.48\textwidth}
\textbf{Đã xác nhận}
\begin{itemize}
\item Không có null ở 12 cột.
\item Không có grain duplicate.
\item DAY và WEEK\_NO liên tục 100\%.
\item Không có HHMM bất khả thi.
\item Tất cả PRODUCT\_ID giao dịch join được product master.
\end{itemize}
\column{.48\textwidth}
\textbf{Chưa nên tự động sửa}
\begin{itemize}
\item QUANTITY max 89.638 và 14.466 dòng bằng 0.
\item SALES\_VALUE max 840 và 18.850 dòng bằng 0.
\item 36 RETAIL\_DISC dương.
\item Coupon discount có IQR bằng 0 nhưng giá trị âm mang ý nghĩa giao dịch.
\item Currency và đơn vị QUANTITY chưa được xác nhận.
\end{itemize}
\end{columns}
\medskip
\alert{Nguyên tắc:} ghi nhận extreme trước, chỉ làm sạch sau khi xác nhận ý nghĩa nghiệp vụ.
\src{Schema validation, domain checks và distribution profiles.}
\end{frame}

\section{Raw data có thể tạo tín hiệu propensity nào}

\begin{frame}{Transaction hỗ trợ RFM, preference và promo sensitivity}
\small
\begin{tabularx}{\textwidth}{@{}p{2.5cm}p{3.2cm}Y@{}}
\toprule
\textbf{Tín hiệu} & \textbf{Raw columns} & \textbf{Feature theo household tại snapshot $t$} \\
\midrule
Recency & \code{DAY}, \code{household_key} & $t-\max(DAY)$, số ngày từ lần mua gần nhất \\
\addlinespace
Frequency & \code{BASKET_ID}, \code{DAY} & Số basket, active days, median interpurchase gap trong observation window \\
\addlinespace
Monetary & \code{SALES_VALUE} & Tổng chi tiêu, average basket value, spend trend \\
\addlinespace
Preference & \code{PRODUCT_ID} + product taxonomy & Tỷ trọng chi tiêu theo department, commodity và brand \\
\addlinespace
Promo sensitivity & Ba discount columns & Discount share, coupon-use rate, spend khi có và không có promotion \\
\addlinespace
Routine & \code{STORE_ID}, \code{TRANS_TIME} & Store concentration, shopping-hour entropy, regularity \\
\bottomrule
\end{tabularx}
\src{Suy luận feature engineering từ grain và semantic role của raw columns.}
\end{frame}

\begin{frame}{Lịch sử đủ dài cho time-window propensity}
\begin{columns}[T]
\column{.57\textwidth}
\centering
\begin{tikzpicture}[x=.68cm,y=1cm]
\fill[blue!18] (0,0) rectangle (10,1);
\fill[orange!35] (10,0) rectangle (12.5,1);
\draw[thick] (10,-.2)--(10,1.2);
\node[align=center] at (5,.5){\textbf{Observation window}\\12 tháng};
\node[align=center] at (11.25,.5){\textbf{Prediction}\\30 ngày};
\node[below,font=\scriptsize] at (0,0){$t-365$};
\node[below,font=\scriptsize] at (10,0){Snapshot $t$};
\node[below,font=\scriptsize] at (12.5,0){$t+30$};
\end{tikzpicture}
\column{.39\textwidth}
\textbf{Bằng chứng khả dụng}
\begin{itemize}
\item Median household span: 627 ngày.
\item 2.461/2.500 households có span ít nhất 365 ngày.
\item Median 79 baskets và 70 active days.
\end{itemize}
\end{columns}
\medskip
\begin{alertblock}{Kiểm soát leakage}
Feature chỉ dùng giao dịch có DAY $\leq t$. Nhãn phải nằm hoàn toàn sau snapshot. Split train, validation và test theo DAY.
\end{alertblock}
\src{Household-level aggregation và nguyên tắc time-window modeling.}
\end{frame}

\begin{frame}{Kết luận và câu hỏi cần chốt}
\begin{columns}[T]
\column{.48\textwidth}
\textbf{Transaction phù hợp làm bảng fact}
\begin{itemize}
\item Grain rõ và unique.
\item Household lặp lại qua thời gian dài.
\item Đủ key để nối product, promotion và demographic.
\item Raw columns hỗ trợ RFM, preference và promo sensitivity.
\end{itemize}
\column{.48\textwidth}
\textbf{Bốn câu hỏi trước khi làm sạch}
\begin{enumerate}
\item DAY 1 tương ứng ngày lịch nào?
\item QUANTITY dùng đơn vị gì?
\item Currency và quy ước dấu discount là gì?
\item Zero quantity hoặc zero sales mang nghĩa nghiệp vụ gì?
\end{enumerate}
\end{columns}
\medskip
\centering
\color{structure.fg}\bfseries Bước tiếp theo: join product taxonomy để giải thích preference signals.
\src{Tổng hợp EDA transaction\_data.csv.}
\end{frame}

\end{document}
